% !PRINTED none
% The source page between printed pp. 521 and 523 is unnumbered; its complete
% section-opening content is translated below.
\begin{center}
\gesperrt{Fifteenth Section.}

\gesperrt{Cyclic Equations.}

\rule{3cm}{0.4pt}

\S. 159.

\gesperrt{Cubic Equations.}
\end{center}

From the standpoint now gained we shall first consider the
solution of the equations of the third and fourth degree.

The group of the permutations of three figures consists of
six elements, namely of the identical permutation 1, of two
three-term cycles, and three transpositions:
\weq{1}{\begin{array}{ccc} 1 & (0,\,1,\,2) & (0,\,2,\,1) \\ (1,\,2) & (2,\,0) & (0,\,1). \end{array}}
The first three
\[
1, \; (0,\,1,\,2) \quad (0,\,2,\,1)
\]
form the alternating group; it consists of the powers of the
cyclic permutation $\pi = (0, 1, 2)$:
\weq{2}{\pi^0 = 1, \; \pi, \; \pi^2}
and is therefore called the \gesperrt{cyclic group}.

Now let $\alpha, \alpha_1, \alpha_2$ be the three roots of the cubic
equation:
\weq{3}{f(x) = x^3 - ax^2 + bx - c = 0.}
If we take as basis the field $\Omega$, which consists of all rational
functions of the independent variables $a, b, c$,
then (1) is the group of this equation; but if we
\weq{4}{\sqrt{D} = (\alpha - \alpha_1)(\alpha - \alpha_2)(\alpha_1 - \alpha_2)}
%CONT>>
